% ============================================================
% TA-IND-01 — Informe Académico Individual
% Análisis crítico de Apache Kafka frente a RabbitMQ
% Compilar: pdflatex -> biber -> pdflatex -> pdflatex
% ============================================================
\documentclass[12pt,a4paper]{article}

\usepackage[utf8]{inputenc}
\usepackage[T1]{fontenc}
\usepackage[spanish,es-tabla]{babel}
\usepackage{csquotes}
\usepackage[margin=2.5cm]{geometry}
\usepackage{graphicx}
\usepackage{float}
\usepackage{caption}
\usepackage{booktabs}
\usepackage{tabularx}
\usepackage{xcolor}
\usepackage{tikz}
\usepackage[hidelinks]{hyperref}
\usepackage[backend=biber,style=ieee]{biblatex}
\addbibresource{referencias.bib}

\graphicspath{{figuras/}{capturas/}{anexos/}}

% ---------- Datos de la carátula ----------
\newcommand{\estudiante}{Rivas Piloso Dayver Yahir}
\newcommand{\paralelo}{Séptimo nivel -- Paralelo B}
\newcommand{\fechaentrega}{15 de noviembre de 2026}
\newcommand{\repo}{https://github.com/drivasp/TA-IND-01-Kafka}

\begin{document}

% ============================================================
% 1. CARÁTULA
% ============================================================
\begin{titlepage}
  \centering
  \IfFileExists{figuras/logo_uteq.jpeg}{\includegraphics[width=0.28\textwidth]{logo_uteq.jpeg}\par\vspace{0.8cm}}{}
  {\Large\bfseries UNIVERSIDAD TÉCNICA ESTATAL DE QUEVEDO\par}
  \vspace{0.4cm}
  {\large Facultad de Ciencias de la Computación\par}
  {\large Carrera de Software\par}
  \vspace{1cm}
  {\large\bfseries Aplicaciones Distribuidas (ISR-701)\par}
  \vspace{0.2cm}
  {\paralelo\par}
  \vspace{1cm}
  {\large\bfseries TA-IND-01 --- Informe Académico Individual\par}
  \vspace{0.2cm}
  {\large Análisis crítico de Apache Kafka frente a RabbitMQ\par}
  \vfill
  \begin{tabular}{rl}
    \textbf{Estudiante:} & \estudiante \\
    \textbf{Docente:}    & Gleiston Cicerón Guerrero Ulloa, PhD \\
    \textbf{Fecha:}      & \fechaentrega \\
    \textbf{Periodo:}    & SPA 2026-2027 \\
    \textbf{Repositorio:} & \url{\repo} \\
  \end{tabular}
  \vspace{1cm}
\end{titlepage}

% ============================================================
% 2. RESUMEN (150 a 200 palabras)
% ============================================================
\begin{abstract}
% Sistema analizado, enfoque y conclusiones principales.
\end{abstract}

\newpage
\tableofcontents
\newpage

% ============================================================
% 3. INTRODUCCIÓN
% ============================================================
\section{Introducción}
% Qué problema resuelve Kafka, quién lo usa, a qué escala y por qué se eligió.

% ============================================================
% 4. ARQUITECTURA DEL SISTEMA
% ============================================================
\section{Arquitectura del sistema}
% Componentes (brokers, tópicos, particiones, réplicas, productores,
% consumidores, controladores KRaft) y cómo se comunican.

% \begin{figure}[H]
%   \centering
%   \includegraphics[width=0.9\textwidth]{arquitectura.png}
%   \caption{Arquitectura de Apache Kafka (elaboración propia).}
%   \label{fig:arquitectura}
% \end{figure}
% Evolución desde el esquema a mano: ver Anexo~\ref{anexo:a}.

% ============================================================
% 5. ANÁLISIS CON LOS CONCEPTOS DE LA UNIDAD 1
% ============================================================
\section{Análisis con los conceptos de la Unidad 1}

\subsection{Transparencias de distribución}
% Acceso, ubicación, migración, reubicación, replicación, concurrencia y fallos.
% Cuáles ofrece y cuáles no, y por qué.

\subsection{Teorema CAP y PACELC}
% Qué prioriza con partición y qué prioriza sin partición. Con evidencia citada.

\subsection{Modelos de comunicación}
% Cliente-servidor, P2P o híbrido; síncrona o asíncrona; sockets, RPC o mensajería.

\subsection{Tolerancia a fallos}
% Qué fallos soporta, cómo los detecta y cómo se recupera. Incidente real documentado.

% ============================================================
% 6. EXPERIMENTACIÓN PROPIA
% ============================================================
\section{Experimentación propia}
% Entorno, comandos (ver experimento/comandos.md), capturas e interpretación.

% \begin{figure}[H]
%   \centering
%   \includegraphics[width=0.9\textwidth]{sistema/captura1.png}
%   \caption{Clúster de tres brokers en ejecución.}
% \end{figure}

% ============================================================
% 7. EJEMPLOS DEL ENTORNO DEL ESTUDIANTE
% ============================================================
\section{Ejemplos del entorno del estudiante}
% Al menos dos sistemas reales usados de primera mano, relacionados con un concepto.

% ============================================================
% 8. ANÁLISIS COMPARATIVO
% ============================================================
\section{Análisis comparativo}

\begin{table}[H]
  \centering
  \caption{Comparación entre Apache Kafka y RabbitMQ}
  \label{tab:comparacion}
  \begin{tabularx}{\textwidth}{@{}lXX@{}}
    \toprule
    \textbf{Criterio} & \textbf{Apache Kafka} & \textbf{RabbitMQ} \\
    \midrule
    Transparencias &  &  \\
    CAP (con partición) &  &  \\
    PACELC (sin partición) &  &  \\
    Modelo de comunicación &  &  \\
    Replicación &  &  \\
    Detección y recuperación de fallos &  &  \\
    \bottomrule
  \end{tabularx}
\end{table}

% Discusión de la tabla.

% ============================================================
% 9. POSTURA CRÍTICA
% ============================================================
\section{Postura crítica}
% Fortalezas, limitaciones, trade-offs y contextos en los que no sería la mejor opción.

% ============================================================
% 10. CONCLUSIONES
% ============================================================
\section{Conclusiones}

% ============================================================
% 11. DECLARACIÓN DE USO DE IA
% ============================================================
\section*{Declaración de uso de inteligencia artificial}
\addcontentsline{toc}{section}{Declaración de uso de inteligencia artificial}
% Herramienta, versión, para qué se usó y cómo se verificó; o declarar que no se usó.

% ============================================================
% 12. REFERENCIAS (IEEE, biblatex)
% ============================================================
\printbibliography[heading=bibintoc,title={Referencias}]

% ============================================================
% 13. ANEXOS
% ============================================================
\appendix
\renewcommand{\thesection}{\Alph{section}}

\section{Esquema inicial a mano}\label{anexo:a}
% \includegraphics[width=0.8\textwidth]{esquema_mano.jpg}

\section{Capturas de lectura de las fuentes}\label{anexo:b}
% \includegraphics[width=0.9\textwidth]{fuentes/fuente1_p5.png}

\end{document}
